\documentclass[12pt]{article}
\usepackage[utf8]{inputenc}
\usepackage[T1]{fontenc}
\usepackage[margin=1in]{geometry}
\usepackage{setspace}
\usepackage{xcolor}
\usepackage[colorlinks=true, urlcolor=blue!60!black, linkcolor=blue!60!black]{hyperref}
\setlength{\parindent}{0em}
\setlength{\parskip}{0.6em}
\renewcommand{\baselinestretch}{1.0}
\begin{document}
\begin{center}
{\LARGE \textit{Kennedy v. Bremerton School District}, 597 U.S. 507 (2022)\par}
\vspace{0.5em}
{\large Case Brief --- Two-Page Edition\par}
\end{center}

\section{Memory Jogger}
A coach's brief, quiet postgame prayer is protected by the Free Exercise and Free Speech Clauses, and \textit{Lemon} gives way to history and tradition (\href{https://www.law.cornell.edu/supremecourt/text/21-418}{597 U.S. 507}).

\section{Detailed Case Facts}
Bremerton High School football coach Joseph Kennedy knelt at midfield after games to pray; players came to join him, and the practice grew into religious talks. In 2015 the District told him to keep on-duty religious expression ``nondemonstrative.'' He stopped the team prayers but, after announcing his plans publicly, prayed alone at midfield after three October games while players sang the fight song. The District, conceding no evidence of direct coercion, put him on leave.

\section{Procedural History}
The district court granted the District summary judgment, and the Ninth Circuit affirmed, calling the prayers government speech justified by Establishment Clause concerns. The Supreme Court reversed.

\section{Judicial Votes}
6--3. Gorsuch wrote for Roberts, Thomas, Alito, Barrett, and (except Part III--B) Kavanaugh; Thomas and Alito concurred; Sotomayor dissented, joined by Breyer and Kagan.

\section{Holding}
A public employee's brief, quiet, personal religious observance during a lull is protected. The Establishment Clause is read by historical practices and understandings, not \textit{Lemon} or endorsement, and absent coercion it does not justify suppression.

\section{Analysis of Opinions}
Gorsuch held the District's policy neither neutral nor generally applicable, and found the prayers private speech under \textit{Pickering}--\textit{Garcetti} because they were not part of Kennedy's duties. The Clauses are complementary; \textit{Lemon} was ``ambitious, abstract, and ahistorical''; and no student was coerced, unlike the captive audiences of \textit{Lee} and \textit{Santa Fe}. Sotomayor said the majority ``misconstrues the facts,'' ignoring years of audible team prayers, and warned that coaches' power over students makes subtle coercion real.

\section{Examples: Future Applications}
\textbf{Same-Side:} a teacher's silent lunchtime prayer at her desk is protected. \textbf{Opposite-Side:} a principal's prayer over the homeroom intercom is unconstitutional. \textbf{Fence-Sitter:} a coach's silent prayer that students routinely join just before a mandatory team critique.

\section{Critique}
Supporters praise the end of \textit{Lemon}, though some originalists find the historical analysis thin. Lupu and Tuttle call the decision ``a sledgehammer to the bedrock of nonestablishment,'' warning that a direct-coercion test ignores the pressures on students.

\section{Key Quotations}
\begin{quote}
``In place of \textit{Lemon} and the endorsement test, this Court has instructed that the Establishment Clause must be interpreted by `reference to historical practices and understandings.' ''
\end{quote}
\noindent --- Justice Gorsuch, Opinion of the Court, 597 U.S. 507.

\end{document}
